\documentclass[12pt,french]{book}
%%%%%%%%%%%%%%%%%%%%%%%%%%%%%%%%%%%%%%%%%%%%%%%%%%%%%%%%%%%%%%%%%%%%%%%%%%%%%%%
%___________________________
%===    Configurations 12.09.2013
%------------------------------------------------------
%packages permettant d'augmenter le nombre de registres de dimension et donc d'éviter les erreurs de compilation dûs aux packages tikz, pstricks and compagnie
\usepackage{etex}
%___________________________
%===    Pour le français
%------------------------------------------------------
\usepackage[utf8x]{inputenc}
\usepackage[T1]{fontenc}
\usepackage{babel}
\FrenchFootnotes

%___________________________
%===    Polices d'écriture
%------------------------------------------------------
%\usepackage{mathpazo}
\usepackage{frcursive} % Pour l'écriture cursive
\usepackage[upright]{fourier}% l'option permet d'avoir les majuscules droites dans les formules mathématiques
\usepackage[scaled=0.875]{helvet}

%___________________________
%===    Les couleurs
%------------------------------------------------------
\usepackage[dvipsnames,table]{xcolor}
%
\newcommand{\rouge}[1]{{\color{red} #1}}
\definecolor{midblue}{rgb}{0.145,0.490,0.882}
\newcommand\MaCouleur{midblue}

%___________________________
%===   Redéfinition des marges par défaut
%------------------------------------------------------
\setlength\paperheight{297mm}
\setlength\paperwidth{210mm}
\setlength{\evensidemargin}{0cm}% Marge gauche sur pages paires
\setlength{\oddsidemargin}{0cm}%{-0.5cm}% Marge gauche sur pages impaires
\setlength{\topmargin}{-2cm}% Marge en haut
\setlength{\headsep}{0.5cm}% Entre le haut de page et le texte
\setlength{\headheight}{0.7cm}% Haut de page
\setlength{\textheight}{25.2cm}% Hauteur de la zone de texte
\setlength{\textwidth}{17cm}% Largeur de la zone de texte

\usepackage{lscape} %permet le format paysage du document
\usepackage{xspace} % création automatique d'espaces dans les commandes
\setlength{\parindent}{0pt}

\usepackage{fancyhdr}
\pagestyle{fancy}
%
\renewcommand{\headrulewidth}{0pt}% pas de trait en entête
\newcommand\RegleEntete[1][0.4pt]{\renewcommand{\headrulewidth}{#1}}%commande pour ajouter un trait horizontal en entête

\newcommand{\entete}[3]{\lhead{#1} \chead{#2} \rhead{#3}}
\newcommand{\pieddepage}[3]{\lfoot{#1} \cfoot{#2} \rfoot{#3}}
%
%\renewcommand{\chaptermark}[1]{\markboth{#1}{}} % enregistre le titre courant du chapitre
    %en-tete droite page [paire] et {impaire}
%\rhead[]{\textbf{\leftmark.}}
    %en-tete gauche page [paire] et {impaire}
%\lhead[\textbf{\chaptername~\thechapter.}]{}


\usepackage{enumerate} %permet la modif de la numérotation et de poursuivre une numérotation en cours avec \begin{enumerate}[resume]
\usepackage{enumitem}
\frenchbsetup{StandardLists=true}%frenchb ne s'occupera pas des listes
\setenumerate[1]{font=\bfseries,label=\arabic*\degres)} % numérotation 1°) 2°) ...
%\setenumerate[2]{font=\itshape,label=(\alph*)} % sous-numérotation (a) (b) ...
\setenumerate[2]{font=\bfseries,label=(\alph*)} % sous-numérotation (a) (b) ...

\usepackage{lastpage} % permet d'afficher le nombre total de pages après DEUX compilations.

%___________________________
%===    Raccourcis classe
%------------------------------------------------------
\newcommand\seconde{2\up{nde}\xspace}
\newcommand\premiere{1\up{ère}\xspace}
\newcommand\terminale{T\up{le}\xspace}
\newcommand\stmg{\bsc{Stmg}}
\newcommand\sti{\bsc{Sti2d}}
\newcommand\bat{BAT 1\xspace}
\newcommand\BAT{BAT 2\xspace}
\newcommand\tesspe{TES Spécialité\xspace}


%___________________________
%===    Réglages et Commandes Maths
%------------------------------------------------------
%redéfinition de fractions, limites, sommes, intégrales, coefficients binomiaux en displaystyle, limites de suites
\usepackage{amssymb,mathtools}
\let\binomOld\binom
\renewcommand{\binom}{\displaystyle\binomOld}
\let\limOld\lim
\renewcommand{\lim}{\displaystyle\limOld}
\newcommand{\limn}{\lim_{n\to +\infty}} %limite lorsque n tend vers + infini
\newcommand{\limm}{\lim_{x\to -\infty}} %limite lorsque x tend vers - infini
\newcommand{\limp}{\lim_{x\to +\infty}} %limite lorsque x tend vers + infini
\newcommand{\limz}{\lim_{x\to 0}} %limite lorsque x tend vers 0
\newcommand{\limzm}{\lim_{\substack{x \to 0\\ x < 0}}} %limite lorsque x tend vers 0-
\newcommand{\limzp}{\lim_{\substack{x \to 0\\ x > 0}}} %limite lorsque x tend vers 0+
\let\sumOld\sum
\renewcommand{\sum}{\displaystyle\sumOld}
\let\intOld\int
\renewcommand{\int}{\displaystyle\intOld}

%\usepackage{yhmath}%permet les arcs de cercles
%\usepackage[euler-digits]{eulervm} %-> police maths
%
\usepackage{amssymb,mathtools}
\usepackage{stmaryrd}%\llbracket et \rrbracket % crochets doubles pour intervalles d'entier
%symbole parallèle avec \sslash

\newcommand{\crochets}[2]{\ensuremath{\llbracket #1 ; #2 \rrbracket}}

\newcommand{\intervalleff}[2]{\left[#1\,;#2\right]}
\newcommand{\intervallefo}[2]{\left[#1\,;#2\right[}
\newcommand{\intervalleof}[2]{\left]#1\,;#2\right]}
\newcommand{\intervalleoo}[2]{\left]#1\,;#2\right[}



\usepackage{bm} % pour l'écriture en gras des formules mathématiques avec \bm

\usepackage{cancel} % pour les simplifications de fractions
\renewcommand\CancelColor{\color{red}}
%\usepackage{siunitx} % écriture de nombres et d'unités
%\sisetup{output-decimal-marker={,},detect-all}
\usepackage[autolanguage,np]{numprint}
%permet les espacement pour les nombres décimaux avec \np{3,12456} en environnement maths ou pas
\DecimalMathComma %supprime l'espace après la virgule dans un nombre

%
\usepackage{dsfont} %écriture des ensemble N, R, C ...
\newcommand{\C}{\mathds C}
\newcommand{\R}{\mathds R}
\newcommand{\Q}{\mathds Q}
\newcommand{\D}{\mathds D}
\newcommand{\Z}{\mathds Z}
\newcommand{\N}{\mathds N}
\newcommand\Ind{\mathds 1} %= fonction indicatrice
\newcommand\p{\mathds P} %= probabilité
\newcommand\E{\mathds E} % Espérance
\newcommand\V{\mathds V} % Variance
\newcommand{\e}{\text{e}}
\newcommand{\dd}{\,\text{d}}

%Nombres complexes
\let\Reold\Re
\renewcommand{\Re}{~\text{Re}~}
\let\Imold\Im
\renewcommand{\Im}{~\text{Im}~}
\newcommand{\ii}{\,\text{i}}
% Exponentielle complexe
\newcommand{\ei}[2]{\,\e^{\dfrac{#1\ii\pi}{#2}}}


%
\usepackage{mathrsfs}   % Police de maths jolie caligraphie
\newcommand{\calig}[1]{\ensuremath{\mathscr{#1}}}
\newcommand\mtc[1]{\ensuremath{\mathcal{#1}}}


%Gestion des espaces
%
\newcommand{\pv}{\ensuremath{\: ; \,}}
\newlength{\EspacePV}
\setlength{\EspacePV}{1em plus 0.5em minus 0.5em}
\newcommand{\qq}{\hspace{\EspacePV} ; \hspace{\EspacePV}}
\newcommand{\qetq}{\hspace{\EspacePV} \text{et} \hspace{\EspacePV}}
\newcommand{\qouq}{\hspace{\EspacePV} \text{ou} \hspace{\EspacePV}}
\newcommand{\qLq}{\hspace{\EspacePV} \Leftarrow \hspace{\EspacePV}}
\newcommand{\qRq}{\hspace{\EspacePV} \Rightarrow \hspace{\EspacePV}}
\newcommand{\qLRq}{\hspace{\EspacePV} \Leftrightarrow \hspace{\EspacePV}}

%simplification notation norme \norme{}
\newcommand{\norme}[1]{\left\Vert #1\right\Vert}


%simplification de la notation de vecteur \vect{}
\newcommand{\vect}[1]{\mathchoice%
{\overrightarrow{\displaystyle\mathstrut#1\,\,}}%
{\overrightarrow{\textstyle\mathstrut#1\,\,}}%
{\overrightarrow{\scriptstyle\mathstrut#1\,\,}}%
{\overrightarrow{\scriptscriptstyle\mathstrut#1\,\,}}}



%Repères
\def\Oij{$\left(\text{O}\pv\vect{\imath},~\vect{\jmath}\right)$\xspace}
\def\Oijk{$\left(\text{O}\pv\vect{\imath},~ \vect{\jmath},~ \vect{k}\right)$\xspace}
\def\Ouv{$\left(\text{O}\pv\vect{u},~\vect{v}\right)$\xspace}
\def\OIJ{$\left(O\pv I\:,\,J\right)$\xspace}

\newcommand\abs[1]{\ensuremath{\left\vert #1 \right\vert}}%valeur absolue
\newcommand\Arc[1]{\ensuremath{\wideparen{#1}}}%arc de cercle


%symbole pour variable aléatoire qui suit une loi
\newcommand{\suit}{\hookrightarrow}

%___________________________
%===    Pour les tableaux
%------------------------------------------------------
\usepackage{array}
\usepackage{longtable}
\usepackage{tabularx,tabulary}
\usepackage{multirow}
\usepackage{multicol}
%exemple
%\begin{multicols}{3}[Titre sur une seule colonne.]
%   3~colonnes équilibrées, 3~colonnes équilibrées, 3~colonnes équilibrées, 3~colonnes équilibrées
%\end{multicols}
%\begin{multicols}{2}[\section{Titre numéroté.}]
%   blabla sur deux colonnes, c'est plus sérieux. C'est le style qui est généralement utilisé pour écrire des articles.
%saut de colonne forcé :
%\columnbreak
%djhskjdhjsq
%sdkksqjhd
%\end{multicols}
%Pour ajouter un titre numéroté qui apparaisse sur toute la largeur de la page, il faut utiliser l'option [\section{Titre.}] juste après \begin{multicols}{nb-col}.
%Remarques :
%Pour qu'une ligne de séparation apparaisse entre les colonnes, il faut utiliser : \setlength{\columnseprule}{1pt}.

%Pour redéfinir la largeur de l'espace inter-colonnes, il faut utiliser \setlength{\columnsep}{30pt}.

%Pour remonter le texte, dans chaque colonne vers le haut : \raggedcolumns qui se tape :\begin{multicols}{2}\raggedcolumns...\columnbreak...\columnbreak\end{multicols}

%Pour supprimer les traits verticaux : \setlength{\columnseprule}{0pt} avant \begin{multicols}{3}...\end{multicols}
\setlength\columnseprule{0.4pt}
\renewcommand{\arraystretch}{1.5}%augmente la hauteur des lignes des tableaux
%colonnes centrées verticalement et horizontalement permettant d'écrire des paragraphes de largeur fixée du type M{3cm}
\newcolumntype{M}[1]{>{\centering\arraybackslash}m{#1}}%cellule centrée horizontalement et verticalement
%\arraybackslash permet de continuer à utiliser \\ pour le changement de ligne

\usepackage{arydshln}% permet des filets horizontaux ou verticaux en pointillés avec
%pour les filets horizontaux \hdashline ou \cdashline qui s'utilisent comme \hline ou \cline
% pour les filets verticaux les deux points :


%___________________________
%===    Divers packages
%------------------------------------------------------
\usepackage{bclogo}
\usepackage{textcomp}
\usepackage{eurosym}%avec \EUR{3,12}
\usepackage{soul} % Pour souligner : \ul
\usepackage{ulem} % Pour souligner double : \uuline
                      % Pour souligner ondulé : \uwave
                      % Pour barrer horizontal : \sout
                      % Pour barrer diagonal : \xout
\usepackage{tikz,tkz-tab,tkz-graph}
\usetikzlibrary{calc,shapes,arrows,plotmarks,lindenmayersystems,decorations,decorations.pathreplacing,patterns}
\usepackage{pstricks,pst-plot,pst-text,pstricks-add,pst-eucl,pst-all}


%INTERLIGNES
\usepackage{setspace}
%s'utilise avec \begin{spacing}{''facteur''}
%   […]
%\end{spacing}

%Pointillés sur toute la ligne
\usepackage{multido}
\newcommand{\Pointilles}[1][1]{%
\multido{}{#1}{\makebox[\textwidth]{\dotfill}\\[1.5\parskip]
}}
%commandes : \Pointilles ou \Pointilles[4] pour 4 lignes


%textes à trous
\newlength\lgtrou
\newcommand*\trou[1]{%
\settowidth\lgtrou{#1}%
\makebox[2\lgtrou]{\dotfill}
\setlength\baselineskip{1.2\baselineskip}}
%Commande à utiliser : \trou{texte qui sera remplacé par des pointillés}

%divers cadres
\usepackage{fancybox} % par exemple \ovalbox{}

%caractères spéciaux avec la commande \ding{230} par exemple
\usepackage{pifont}

%___________________________
%===    Quelques raccourcis perso
%------------------------------------------------------
\newcommand\pfr[1]{\psframebox[linecolor=red]{#1}}
\newcommand\coef[1][]{c{\oe}fficient#1\xspace}


%QRcode, codebarre
\usepackage{pst-barcode}
%\begin{pspicture}(2,2)
%	\psbarcode{http://www.latex-howto.be}{eclevel=M}{qrcode}
%\end{pspicture}


%Texte en filigrane
\usepackage{watermark}
%On utilise ensuite les commandes \watermark, \leftwatermark, \rightwatermark ou \thiswatermark qui permettent de définir un filigrane sur toutes les pages, les pages paires, les pages impaires ou juste une page
%Exemple : \thiswatermark {
%\begin{minipage}{0.95\linewidth}
%\vspace{25cm}
%\begin{center}
%\rotatebox{55}{\scalebox{8}{\color[gray]{0.7}\LaTeX}}
%\end{center}
%\end{minipage}
%}

%QCM
\usepackage{alterqcm}					%%Permet de créer des QCM
%\begin{alterqcm}
%\AQquestion{Question}{{Proposition 1},{Proposition 2},{Proposition 3}}
%\end{alterqcm}

%\dingsquare %carré avant V ou F
%\dingchecksquare %carré validé devant V ou F


%Rond entourant une lettre avec pour arguments la couleur de fond, puis la lettre
\newcommand\rond[2][red!20]{\tikz[baseline]{\node[fill=#1,anchor=base,circle]{\bf #2};}}


%Ecrire card en écriture normale :
\newcommand{\card}{\text{card}\xspace}


%___________________________
%===    ALGORITHMES
%------------------------------------------------------

%ALGORITHME avec Algobox
\usepackage{ucs}
\usepackage{framed}
\definecolor{fond}{gray}{0.95}
\newenvironment{cadrecode}{%
  \def\FrameCommand{{\color[HTML]{888888}\vrule width 3pt}\colorbox{fond}}%
  \MakeFramed {\advance\hsize-\width \FrameRestore}}%
{\endMakeFramed}
\usepackage{alltt}

% Mise en forme des algorithmes
\usepackage[french,boxed,titlenumbered,lined,longend]{algorithm2e}
  \SetKwIF {Si}{SinonSi}{Sinon}{si}{alors}{sinon\_si}{alors}{fin~si}
 \SetKwFor{Tq}{tant\_que~}{~faire~}{fin~tant\_que}
 \SetKwFor{PourCh}{pour\_chaque }{ faire }{fin pour\_chaque}
 \SetKwInput{Sortie}{Sortie}
  \SetKwInput{Entree}{Entrée}
\newcommand{\Algocmd}[1]{\textsf{\textsc{\textbf{#1}}}}\SetKwSty{Algocmd}
  \newcommand{\AlgCommentaire}[1]{\textsl{\small  #1}}


%___________________________
%===    MISE EN FORME EXERCICES
%------------------------------------------------------
%\usepackage{marvosym}
\usepackage{slashbox}

\newcounter{exo}
\newenvironment{exo}{%
  \refstepcounter{exo}\Writinghand\ \textbf{Exercice \theexo.}\par
  \medskip}%
{\[*\]}


%___________________________
%===    HYPERLIENS
%------------------------------------------------------
\usepackage[colorlinks=true,linkcolor=black,filecolor=blue,urlcolor=blue,bookmarksnumbered]{hyperref} 

\setlength{\oddsidemargin}{-0.5cm}% Marge gauche sur pages impaires


% Environnement enumerate
\renewcommand{\theenumi}{\bf\textsf{\arabic{enumi}}}
\renewcommand{\labelenumi}{\bf\textsf{\theenumi.}}
\renewcommand{\theenumii}{\bf\textsf{\alph{enumii}}}
\renewcommand{\labelenumii}{\bf\textsf{\theenumii.}}
\renewcommand{\theenumiii}{\bf\textsf{\roman{enumiii}}}
\renewcommand{\labelenumiii}{\bf\textsf{\theenumiii.}}



\usetikzlibrary{shadows,trees}


%definition des couleurs
\definecolor{fondpaille}{cmyk}{0,0,0.1,0}%\pagecolor{fondpaille}
\definecolor{gris}{rgb}{0.7,0.7,0.7}
\definecolor{rouge}{rgb}{1,0,0}
\definecolor{bleu}{rgb}{0,0,1}
\definecolor{vert}{rgb}{0,1,0}
\definecolor{deficolor}{HTML}{2D9AFF}
\definecolor{backdeficolor}{HTML}{EDEDED}%{036DD0}%dégradé bleu{666666}%dégradé gris
\definecolor{theocolor}{HTML}{036DD0}%F4404D%rouge
\definecolor{backtheocolor}{HTML}{D3D3D3}
\definecolor{methcolor}{HTML}{008800}%12BB05}
\definecolor{backmethcolor}{HTML}{FFFACD}
\definecolor{backilluscolor}{HTML}{EDEDED}
\definecolor{sectioncolor}{HTML}{B2B2B2}%vert : {HTML}{008800}%{HTML}{2D9AFF}
\definecolor{subsectioncolor}{HTML}{B2B2B2}%vert : {HTML}{008800}%{rgb}{0.5,0,0}
\definecolor{engcolor}{HTML}{D4D7FE}
\definecolor{exocolor}{rgb}{0,0.6,0}
\definecolor{exosoltitlecolor}{rgb}{0,0.6,0}
\definecolor{titlecolor}{rgb}{1,1,1}

%commande pour enlever les couleurs avant impression
\newcommand{\nocolor}
{\pagecolor{white}
\definecolor{gris}{rgb}{0.7,0.7,0.7}
\definecolor{rouge}{rgb}{0,0,0}
\definecolor{bleu}{rgb}{0,0,0}
\definecolor{vert}{rgb}{0,0,0}
\definecolor{deficolor}{HTML}{B2B2B2}
\definecolor{backdeficolor}{HTML}{EEEEEE}%{036DD0}%dégradé bleu{666666}%dégradé gris
\definecolor{theocolor}{HTML}{B2B2B2}
\definecolor{backtheocolor}{HTML}{EEEEEE}
\definecolor{methcolor}{HTML}{B2B2B2}
\definecolor{backmethcolor}{HTML}{EEEEEE}
\definecolor{backilluscolor}{HTML}{EEEEEE}
\definecolor{sectioncolor}{HTML}{B2B2B2}
\definecolor{subsectioncolor}{HTML}{B2B2B2}
\definecolor{engcolor}{HTML}{EEEEEE}
\definecolor{exocolor}{HTML}{3B3838}
\definecolor{exosoltitlecolor}{rgb}{0,0,0}
\definecolor{titlecolor}{rgb}{0,0,0}
}




%Exercices résolus dans le cours
%#1 : énoncé
%#2 : solution
\newcounter{exosol}
\newcommand{\exosol}[2]{
\stepcounter{exosol}
\begin{tikzpicture}[node distance=0 cm]
\node[fill=backilluscolor,rounded corners=2pt,anchor=south west] (illus) at (0,-0.02)
{\it \textbf{\textcolor{exosoltitlecolor}{Exercice résolu \arabic{exosol}~:~}}};
\node[fill=backilluscolor,rounded corners=2pt,anchor=north west]at(0,0)
{\parbox{\columnwidth-10pt}{#1\par\medskip{\it \textbf{\textcolor{exosoltitlecolor}{Solution~:~}}}#2 }};
\end{tikzpicture}
\bigskip
}

\newcommand{\suite}[1]{
\begin{tikzpicture}[node distance=0 cm]
\node[fill=backilluscolor,rounded corners=2pt,anchor=north west]at(0,0)
{\parbox{\columnwidth-10pt}{{\it \textbf{\textcolor{exosoltitlecolor}{Suite de la solution~:}}}\par#1}};
\end{tikzpicture}
\bigskip
}



%%%%%%%%%%%%%%%%%%%%%%%%%%%%%%%%%%%%%%%%%%%%%%%%%%%%%%%%%%%%%%%%%%%%%%%%%%%%%%%
%Encadrés pour Propriétés, Théorème, Définitions, exemples, exercices

\usepackage{environ}

%\newcounter{propr}
%\newenvironment{propr}{\refstepcounter{propr}\begin{bclogo}[%noborder=true , 
%couleur = white , arrondi = 0.1 ,logo = \bccrayon , barre = snake , tailleOndu = 1.5 , ombre = true , epOmbre = 0.15 , couleurOmbre = black!30]{Propriété \thepropr}}{\end{bclogo}}

\NewEnviron{Prop}[1][]{
\begin{tikzpicture}[node distance=0 cm]
\node[fill=theocolor,rounded corners=5pt,anchor=south west] (theorem) at (0,0)
{\textcolor{titlecolor}{Propriété~:~#1}};
\node[draw,drop shadow,color=theocolor,very thick,fill=backtheocolor,rounded corners=5pt,anchor=north west] at(0,-0.02)
{\black\parbox{\columnwidth-12pt}{\BODY}};
\end{tikzpicture}
\bigskip
}

%\newcommand{\propr}[2]{
%\begin{tikzpicture}[node distance=0 cm]
%\node[fill=theocolor,rounded corners=5pt,anchor=south west] (theorem) at (0,0)
%{\textcolor{titlecolor}{Propriété~:~#1}};
%\node[draw,drop shadow,color=theocolor,very thick,fill=backtheocolor,rounded corners=5pt,anchor=north west] at(0,-0.02)
%{\black\parbox{\columnwidth-12pt}{#2}};
%\end{tikzpicture}
%\bigskip
%}

%\newenvironment{propr}[1][]{\begin{bclogo}[noborder=true , 
%couleur = white , arrondi = 0.1 , barre = none , ombre = true , epOmbre = 0.15 , couleurOmbre = black!30]{Propriété {\bf \xspace : #1}}}{\end{bclogo}}
%
%
%\newcounter{theo}
%\newenvironment{theo}{\refstepcounter{theo}\begin{bclogo}[%noborder=true , 
%couleur = white , arrondi = 0.1 ,logo = \bcbook , barre = snake , tailleOndu = 1.5 , ombre = true , epOmbre = 0.15 , couleurOmbre = orange]{Théorème \thetheo}}{\end{bclogo}}

%Théorème
\NewEnviron{Thm}[1][]{
\begin{tikzpicture}[node distance=0 cm]
\node[fill=theocolor,rounded corners=5pt,anchor=south west] (theorem) at (0,0)
{\textcolor{titlecolor}{Théorème~:~#1}};
\node[draw,drop shadow,color=theocolor,very thick,fill=backtheocolor,rounded corners=5pt,anchor=north west] at(0,-0.02)
{\black\parbox{\columnwidth-12pt}{\BODY}};
\end{tikzpicture}
\bigskip
}



%\newcounter{defi}
%\newenvironment{defi}{\refstepcounter{defi}\begin{bclogo}[%noborder=true , 
%couleur = white , arrondi = 0.1 ,logo = \bcinfo , barre = snake , tailleOndu = 1.5 , ombre = true , epOmbre = 0.15 , couleurOmbre = blue]{Définition \thedefi}}{\end{bclogo}}
\NewEnviron{Defi}[1][]{
\begin{tikzpicture}[node distance=0 cm]
\node[fill=theocolor,rounded corners=5pt,anchor=south west] (theorem) at (0,0)
{\textcolor{titlecolor}{Définition~:~#1}};
\node[draw,drop shadow,color=deficolor,very thick,fill=backdeficolor,rounded corners=5pt,anchor=north west] at(0,-0.02)
{\black\parbox{\columnwidth-12pt}{\BODY}};
\end{tikzpicture}
\bigskip
}

%\newcounter{methode}
%\newenvironment{methode}{\refstepcounter{methode}\begin{bclogo}[%noborder=true , 
%couleur = white , arrondi = 0.1 ,logo = \bccrayon , barre = snake , tailleOndu = 1.5 , ombre = true , epOmbre = 0.15 , couleurOmbre = red]{Méthode \themethode}}{\end{bclogo}}
\NewEnviron{Methode}[1][]{
\begin{tikzpicture}[node distance=0 cm]
\node[fill=theocolor,rounded corners=5pt,anchor=south west] (theorem) at (0,0)
{\textcolor{titlecolor}{Méthode~:~#1}};
\node[draw,drop shadow,color=methcolor,very thick,fill=backmethcolor,rounded corners=5pt,anchor=north west] at(0,-0.02)
{\black\parbox{\columnwidth-12pt}{\BODY}};
\end{tikzpicture}
\bigskip
}


%chapitres
\makeatletter

\renewcommand{\@makechapterhead}[1]{
\begin{tikzpicture}
\node[fill=theocolor,rectangle,rounded corners=5pt]{%
\begin{minipage}{\linewidth}
\begin{center}
\vspace*{9pt}
\textcolor{titlecolor}{\Large \textsc{\textbf{Chapitre \thechapter \ : \ #1}}}
\vspace*{9pt}
\end{center}
\end{minipage}
};\end{tikzpicture}
}

\makeatother

\NewEnviron{Exemple}[1][]{
\begin{tikzpicture}[node distance=0 cm]
\node[draw,drop shadow,color=methcolor,very thick,fill=backmethcolor,rounded corners=5pt,anchor=north west] at(0,-0.02)
{\black\parbox{\columnwidth-12pt}{\textbf{Exemple~:~#1}\\
\BODY}};
\end{tikzpicture}
\bigskip
}

\newcounter{exemple}\newcommand{\exemple}{\refstepcounter{exemple}\textbf{Exemple \theexemple \ :}\xspace}
\newcounter{exercice}\newcommand{\exercice}{\refstepcounter{exercice}\textbf{\large{Exercice \theexercice \ :}}\xspace}

\newcounter{probleme}\newcommand{\probleme}{\refstepcounter{probleme}\textbf{\large{Problème \theprobleme \ :}}\xspace}
\newcounter{remarque}\newcommand{\remarque}{\refstepcounter{remarque}\textbf{Remarque \theremarque \ :}\xspace}

\NewEnviron{Rmq}{
\textbf{\large{Remarque :}}\par
\BODY
\bigskip
}

\newcounter{rem}\newcommand{\rem}{\refstepcounter{rem}\textbf{R \therem \ :}\xspace}

%Exercices sans numérotation automatique
%\newcommand{\Exercice}[1]{\textbf{\large{Exercice #1 :}\xspace}}
%Exercices du contrôle numérotés
\newcounter{exercice}
\NewEnviron{Exercice}{
\refstepcounter{exercice}\textbf{\large{Exercice \theexercice \ :}}\par
\BODY
\bigskip
}

%Exercices non numérotés
\NewEnviron{Exo}[1][]{
\textbf{\large{Exercice #1 \ :}}\par
\BODY
\bigskip
}


%\Leftrightarrow
\newcommand{\Lr}{\Leftrightarrow}

\pagecolor{white}

%Sommaire dans les chapitres
\usepackage{minitoc}

%%%%%%%%%%%%%%%%%%%%%%%%%%%%%%%%%%%%%%%%%%%%%%%%%%%%%%%%%%%%%%%%%%%%%%%%%%%%%%%
%%___________________________
%===    Configurations 12.09.2013
%------------------------------------------------------
%packages permettant d'augmenter le nombre de registres de dimension et donc d'éviter les erreurs de compilation dûs aux packages tikz, pstricks and compagnie
\usepackage{etex}
%___________________________
%===    Pour le français
%------------------------------------------------------
\usepackage[utf8x]{inputenc}
\usepackage[T1]{fontenc}
\usepackage{babel}
\FrenchFootnotes

%___________________________
%===    Polices d'écriture
%------------------------------------------------------
%\usepackage{mathpazo}
\usepackage{frcursive} % Pour l'écriture cursive
\usepackage[upright]{fourier}% l'option permet d'avoir les majuscules droites dans les formules mathématiques
\usepackage[scaled=0.875]{helvet}

%___________________________
%===    Les couleurs
%------------------------------------------------------
\usepackage[dvipsnames,table]{xcolor}
%
\newcommand{\rouge}[1]{{\color{red} #1}}
\definecolor{midblue}{rgb}{0.145,0.490,0.882}
\newcommand\MaCouleur{midblue}

%___________________________
%===   Redéfinition des marges par défaut
%------------------------------------------------------
\setlength\paperheight{297mm}
\setlength\paperwidth{210mm}
\setlength{\evensidemargin}{0cm}% Marge gauche sur pages paires
\setlength{\oddsidemargin}{0cm}%{-0.5cm}% Marge gauche sur pages impaires
\setlength{\topmargin}{-2cm}% Marge en haut
\setlength{\headsep}{0.5cm}% Entre le haut de page et le texte
\setlength{\headheight}{0.7cm}% Haut de page
\setlength{\textheight}{25.2cm}% Hauteur de la zone de texte
\setlength{\textwidth}{17cm}% Largeur de la zone de texte

\usepackage{lscape} %permet le format paysage du document
\usepackage{xspace} % création automatique d'espaces dans les commandes
\setlength{\parindent}{0pt}

\usepackage{fancyhdr}
\pagestyle{fancy}
%
\renewcommand{\headrulewidth}{0pt}% pas de trait en entête
\newcommand\RegleEntete[1][0.4pt]{\renewcommand{\headrulewidth}{#1}}%commande pour ajouter un trait horizontal en entête

\newcommand{\entete}[3]{\lhead{#1} \chead{#2} \rhead{#3}}
\newcommand{\pieddepage}[3]{\lfoot{#1} \cfoot{#2} \rfoot{#3}}
%
%\renewcommand{\chaptermark}[1]{\markboth{#1}{}} % enregistre le titre courant du chapitre
    %en-tete droite page [paire] et {impaire}
%\rhead[]{\textbf{\leftmark.}}
    %en-tete gauche page [paire] et {impaire}
%\lhead[\textbf{\chaptername~\thechapter.}]{}


\usepackage{enumerate} %permet la modif de la numérotation et de poursuivre une numérotation en cours avec \begin{enumerate}[resume]
\usepackage{enumitem}
\frenchbsetup{StandardLists=true}%frenchb ne s'occupera pas des listes
\setenumerate[1]{font=\bfseries,label=\arabic*\degres)} % numérotation 1°) 2°) ...
%\setenumerate[2]{font=\itshape,label=(\alph*)} % sous-numérotation (a) (b) ...
\setenumerate[2]{font=\bfseries,label=(\alph*)} % sous-numérotation (a) (b) ...

\usepackage{lastpage} % permet d'afficher le nombre total de pages après DEUX compilations.

%___________________________
%===    Raccourcis classe
%------------------------------------------------------
\newcommand\seconde{2\up{nde}\xspace}
\newcommand\premiere{1\up{ère}\xspace}
\newcommand\terminale{T\up{le}\xspace}
\newcommand\stmg{\bsc{Stmg}}
\newcommand\sti{\bsc{Sti2d}}
\newcommand\bat{BAT 1\xspace}
\newcommand\BAT{BAT 2\xspace}
\newcommand\tesspe{TES Spécialité\xspace}


%___________________________
%===    Réglages et Commandes Maths
%------------------------------------------------------
%redéfinition de fractions, limites, sommes, intégrales, coefficients binomiaux en displaystyle, limites de suites
\usepackage{amssymb,mathtools}
\let\binomOld\binom
\renewcommand{\binom}{\displaystyle\binomOld}
\let\limOld\lim
\renewcommand{\lim}{\displaystyle\limOld}
\newcommand{\limn}{\lim_{n\to +\infty}} %limite lorsque n tend vers + infini
\newcommand{\limm}{\lim_{x\to -\infty}} %limite lorsque x tend vers - infini
\newcommand{\limp}{\lim_{x\to +\infty}} %limite lorsque x tend vers + infini
\newcommand{\limz}{\lim_{x\to 0}} %limite lorsque x tend vers 0
\newcommand{\limzm}{\lim_{\substack{x \to 0\\ x < 0}}} %limite lorsque x tend vers 0-
\newcommand{\limzp}{\lim_{\substack{x \to 0\\ x > 0}}} %limite lorsque x tend vers 0+
\let\sumOld\sum
\renewcommand{\sum}{\displaystyle\sumOld}
\let\intOld\int
\renewcommand{\int}{\displaystyle\intOld}

%\usepackage{yhmath}%permet les arcs de cercles
%\usepackage[euler-digits]{eulervm} %-> police maths
%
\usepackage{amssymb,mathtools}
\usepackage{stmaryrd}%\llbracket et \rrbracket % crochets doubles pour intervalles d'entier
%symbole parallèle avec \sslash

\newcommand{\crochets}[2]{\ensuremath{\llbracket #1 ; #2 \rrbracket}}

\newcommand{\intervalleff}[2]{\left[#1\,;#2\right]}
\newcommand{\intervallefo}[2]{\left[#1\,;#2\right[}
\newcommand{\intervalleof}[2]{\left]#1\,;#2\right]}
\newcommand{\intervalleoo}[2]{\left]#1\,;#2\right[}



\usepackage{bm} % pour l'écriture en gras des formules mathématiques avec \bm

\usepackage{cancel} % pour les simplifications de fractions
\renewcommand\CancelColor{\color{red}}
%\usepackage{siunitx} % écriture de nombres et d'unités
%\sisetup{output-decimal-marker={,},detect-all}
\usepackage[autolanguage,np]{numprint}
%permet les espacement pour les nombres décimaux avec \np{3,12456} en environnement maths ou pas
\DecimalMathComma %supprime l'espace après la virgule dans un nombre

%
\usepackage{dsfont} %écriture des ensemble N, R, C ...
\newcommand{\C}{\mathds C}
\newcommand{\R}{\mathds R}
\newcommand{\Q}{\mathds Q}
\newcommand{\D}{\mathds D}
\newcommand{\Z}{\mathds Z}
\newcommand{\N}{\mathds N}
\newcommand\Ind{\mathds 1} %= fonction indicatrice
\newcommand\p{\mathds P} %= probabilité
\newcommand\E{\mathds E} % Espérance
\newcommand\V{\mathds V} % Variance
\newcommand{\e}{\text{e}}
\newcommand{\dd}{\,\text{d}}

%Nombres complexes
\let\Reold\Re
\renewcommand{\Re}{~\text{Re}~}
\let\Imold\Im
\renewcommand{\Im}{~\text{Im}~}
\newcommand{\ii}{\,\text{i}}
% Exponentielle complexe
\newcommand{\ei}[2]{\,\e^{\dfrac{#1\ii\pi}{#2}}}


%
\usepackage{mathrsfs}   % Police de maths jolie caligraphie
\newcommand{\calig}[1]{\ensuremath{\mathscr{#1}}}
\newcommand\mtc[1]{\ensuremath{\mathcal{#1}}}


%Gestion des espaces
%
\newcommand{\pv}{\ensuremath{\: ; \,}}
\newlength{\EspacePV}
\setlength{\EspacePV}{1em plus 0.5em minus 0.5em}
\newcommand{\qq}{\hspace{\EspacePV} ; \hspace{\EspacePV}}
\newcommand{\qetq}{\hspace{\EspacePV} \text{et} \hspace{\EspacePV}}
\newcommand{\qouq}{\hspace{\EspacePV} \text{ou} \hspace{\EspacePV}}
\newcommand{\qLq}{\hspace{\EspacePV} \Leftarrow \hspace{\EspacePV}}
\newcommand{\qRq}{\hspace{\EspacePV} \Rightarrow \hspace{\EspacePV}}
\newcommand{\qLRq}{\hspace{\EspacePV} \Leftrightarrow \hspace{\EspacePV}}

%simplification notation norme \norme{}
\newcommand{\norme}[1]{\left\Vert #1\right\Vert}


%simplification de la notation de vecteur \vect{}
\newcommand{\vect}[1]{\mathchoice%
{\overrightarrow{\displaystyle\mathstrut#1\,\,}}%
{\overrightarrow{\textstyle\mathstrut#1\,\,}}%
{\overrightarrow{\scriptstyle\mathstrut#1\,\,}}%
{\overrightarrow{\scriptscriptstyle\mathstrut#1\,\,}}}



%Repères
\def\Oij{$\left(\text{O}\pv\vect{\imath},~\vect{\jmath}\right)$\xspace}
\def\Oijk{$\left(\text{O}\pv\vect{\imath},~ \vect{\jmath},~ \vect{k}\right)$\xspace}
\def\Ouv{$\left(\text{O}\pv\vect{u},~\vect{v}\right)$\xspace}
\def\OIJ{$\left(O\pv I\:,\,J\right)$\xspace}

\newcommand\abs[1]{\ensuremath{\left\vert #1 \right\vert}}%valeur absolue
\newcommand\Arc[1]{\ensuremath{\wideparen{#1}}}%arc de cercle


%symbole pour variable aléatoire qui suit une loi
\newcommand{\suit}{\hookrightarrow}

%___________________________
%===    Pour les tableaux
%------------------------------------------------------
\usepackage{array}
\usepackage{longtable}
\usepackage{tabularx,tabulary}
\usepackage{multirow}
\usepackage{multicol}
%exemple
%\begin{multicols}{3}[Titre sur une seule colonne.]
%   3~colonnes équilibrées, 3~colonnes équilibrées, 3~colonnes équilibrées, 3~colonnes équilibrées
%\end{multicols}
%\begin{multicols}{2}[\section{Titre numéroté.}]
%   blabla sur deux colonnes, c'est plus sérieux. C'est le style qui est généralement utilisé pour écrire des articles.
%saut de colonne forcé :
%\columnbreak
%djhskjdhjsq
%sdkksqjhd
%\end{multicols}
%Pour ajouter un titre numéroté qui apparaisse sur toute la largeur de la page, il faut utiliser l'option [\section{Titre.}] juste après \begin{multicols}{nb-col}.
%Remarques :
%Pour qu'une ligne de séparation apparaisse entre les colonnes, il faut utiliser : \setlength{\columnseprule}{1pt}.

%Pour redéfinir la largeur de l'espace inter-colonnes, il faut utiliser \setlength{\columnsep}{30pt}.

%Pour remonter le texte, dans chaque colonne vers le haut : \raggedcolumns qui se tape :\begin{multicols}{2}\raggedcolumns...\columnbreak...\columnbreak\end{multicols}

%Pour supprimer les traits verticaux : \setlength{\columnseprule}{0pt} avant \begin{multicols}{3}...\end{multicols}
\setlength\columnseprule{0.4pt}
\renewcommand{\arraystretch}{1.5}%augmente la hauteur des lignes des tableaux
%colonnes centrées verticalement et horizontalement permettant d'écrire des paragraphes de largeur fixée du type M{3cm}
\newcolumntype{M}[1]{>{\centering\arraybackslash}m{#1}}%cellule centrée horizontalement et verticalement
%\arraybackslash permet de continuer à utiliser \\ pour le changement de ligne

\usepackage{arydshln}% permet des filets horizontaux ou verticaux en pointillés avec
%pour les filets horizontaux \hdashline ou \cdashline qui s'utilisent comme \hline ou \cline
% pour les filets verticaux les deux points :


%___________________________
%===    Divers packages
%------------------------------------------------------
\usepackage{bclogo}
\usepackage{textcomp}
\usepackage{eurosym}%avec \EUR{3,12}
\usepackage{soul} % Pour souligner : \ul
\usepackage{ulem} % Pour souligner double : \uuline
                      % Pour souligner ondulé : \uwave
                      % Pour barrer horizontal : \sout
                      % Pour barrer diagonal : \xout
\usepackage{tikz,tkz-tab,tkz-graph}
\usetikzlibrary{calc,shapes,arrows,plotmarks,lindenmayersystems,decorations,decorations.pathreplacing,patterns}
\usepackage{pstricks,pst-plot,pst-text,pstricks-add,pst-eucl,pst-all}


%INTERLIGNES
\usepackage{setspace}
%s'utilise avec \begin{spacing}{''facteur''}
%   […]
%\end{spacing}

%Pointillés sur toute la ligne
\usepackage{multido}
\newcommand{\Pointilles}[1][1]{%
\multido{}{#1}{\makebox[\textwidth]{\dotfill}\\[1.5\parskip]
}}
%commandes : \Pointilles ou \Pointilles[4] pour 4 lignes


%textes à trous
\newlength\lgtrou
\newcommand*\trou[1]{%
\settowidth\lgtrou{#1}%
\makebox[2\lgtrou]{\dotfill}
\setlength\baselineskip{1.2\baselineskip}}
%Commande à utiliser : \trou{texte qui sera remplacé par des pointillés}

%divers cadres
\usepackage{fancybox} % par exemple \ovalbox{}

%caractères spéciaux avec la commande \ding{230} par exemple
\usepackage{pifont}

%___________________________
%===    Quelques raccourcis perso
%------------------------------------------------------
\newcommand\pfr[1]{\psframebox[linecolor=red]{#1}}
\newcommand\coef[1][]{c{\oe}fficient#1\xspace}


%QRcode, codebarre
\usepackage{pst-barcode}
%\begin{pspicture}(2,2)
%	\psbarcode{http://www.latex-howto.be}{eclevel=M}{qrcode}
%\end{pspicture}


%Texte en filigrane
\usepackage{watermark}
%On utilise ensuite les commandes \watermark, \leftwatermark, \rightwatermark ou \thiswatermark qui permettent de définir un filigrane sur toutes les pages, les pages paires, les pages impaires ou juste une page
%Exemple : \thiswatermark {
%\begin{minipage}{0.95\linewidth}
%\vspace{25cm}
%\begin{center}
%\rotatebox{55}{\scalebox{8}{\color[gray]{0.7}\LaTeX}}
%\end{center}
%\end{minipage}
%}

%QCM
\usepackage{alterqcm}					%%Permet de créer des QCM
%\begin{alterqcm}
%\AQquestion{Question}{{Proposition 1},{Proposition 2},{Proposition 3}}
%\end{alterqcm}

%\dingsquare %carré avant V ou F
%\dingchecksquare %carré validé devant V ou F


%Rond entourant une lettre avec pour arguments la couleur de fond, puis la lettre
\newcommand\rond[2][red!20]{\tikz[baseline]{\node[fill=#1,anchor=base,circle]{\bf #2};}}


%Ecrire card en écriture normale :
\newcommand{\card}{\text{card}\xspace}


%___________________________
%===    ALGORITHMES
%------------------------------------------------------

%ALGORITHME avec Algobox
\usepackage{ucs}
\usepackage{framed}
\definecolor{fond}{gray}{0.95}
\newenvironment{cadrecode}{%
  \def\FrameCommand{{\color[HTML]{888888}\vrule width 3pt}\colorbox{fond}}%
  \MakeFramed {\advance\hsize-\width \FrameRestore}}%
{\endMakeFramed}
\usepackage{alltt}

% Mise en forme des algorithmes
\usepackage[french,boxed,titlenumbered,lined,longend]{algorithm2e}
  \SetKwIF {Si}{SinonSi}{Sinon}{si}{alors}{sinon\_si}{alors}{fin~si}
 \SetKwFor{Tq}{tant\_que~}{~faire~}{fin~tant\_que}
 \SetKwFor{PourCh}{pour\_chaque }{ faire }{fin pour\_chaque}
 \SetKwInput{Sortie}{Sortie}
  \SetKwInput{Entree}{Entrée}
\newcommand{\Algocmd}[1]{\textsf{\textsc{\textbf{#1}}}}\SetKwSty{Algocmd}
  \newcommand{\AlgCommentaire}[1]{\textsl{\small  #1}}


%___________________________
%===    MISE EN FORME EXERCICES
%------------------------------------------------------
%\usepackage{marvosym}
\usepackage{slashbox}

\newcounter{exo}
\newenvironment{exo}{%
  \refstepcounter{exo}\Writinghand\ \textbf{Exercice \theexo.}\par
  \medskip}%
{\[*\]}


%___________________________
%===    HYPERLIENS
%------------------------------------------------------
\usepackage[colorlinks=true,linkcolor=black,filecolor=blue,urlcolor=blue,bookmarksnumbered]{hyperref} 

%\setlength{\oddsidemargin}{-0.5cm}% Marge gauche sur pages impaires


% Environnement enumerate
\renewcommand{\theenumi}{\bf\textsf{\arabic{enumi}}}
\renewcommand{\labelenumi}{\bf\textsf{\theenumi.}}
\renewcommand{\theenumii}{\bf\textsf{\alph{enumii}}}
\renewcommand{\labelenumii}{\bf\textsf{\theenumii.}}
\renewcommand{\theenumiii}{\bf\textsf{\roman{enumiii}}}
\renewcommand{\labelenumiii}{\bf\textsf{\theenumiii.}}



\usetikzlibrary{shadows,trees}


%definition des couleurs
\definecolor{fondpaille}{cmyk}{0,0,0.1,0}%\pagecolor{fondpaille}
\definecolor{gris}{rgb}{0.7,0.7,0.7}
\definecolor{rouge}{rgb}{1,0,0}
\definecolor{bleu}{rgb}{0,0,1}
\definecolor{vert}{rgb}{0,1,0}
\definecolor{deficolor}{HTML}{2D9AFF}
\definecolor{backdeficolor}{HTML}{EDEDED}%{036DD0}%dégradé bleu{666666}%dégradé gris
\definecolor{theocolor}{HTML}{036DD0}%F4404D%rouge
\definecolor{backtheocolor}{HTML}{D3D3D3}
\definecolor{methcolor}{HTML}{008800}%12BB05}
\definecolor{backmethcolor}{HTML}{FFFACD}
\definecolor{backilluscolor}{HTML}{EDEDED}
\definecolor{sectioncolor}{HTML}{B2B2B2}%vert : {HTML}{008800}%{HTML}{2D9AFF}
\definecolor{subsectioncolor}{HTML}{B2B2B2}%vert : {HTML}{008800}%{rgb}{0.5,0,0}
\definecolor{engcolor}{HTML}{D4D7FE}
\definecolor{exocolor}{rgb}{0,0.6,0}
\definecolor{exosoltitlecolor}{rgb}{0,0.6,0}
\definecolor{titlecolor}{rgb}{1,1,1}

%commande pour enlever les couleurs avant impression
\newcommand{\nocolor}
{\pagecolor{white}
\definecolor{gris}{rgb}{0.7,0.7,0.7}
\definecolor{rouge}{rgb}{0,0,0}
\definecolor{bleu}{rgb}{0,0,0}
\definecolor{vert}{rgb}{0,0,0}
\definecolor{deficolor}{HTML}{B2B2B2}
\definecolor{backdeficolor}{HTML}{EEEEEE}%{036DD0}%dégradé bleu{666666}%dégradé gris
\definecolor{theocolor}{HTML}{B2B2B2}
\definecolor{backtheocolor}{HTML}{EEEEEE}
\definecolor{methcolor}{HTML}{B2B2B2}
\definecolor{backmethcolor}{HTML}{EEEEEE}
\definecolor{backilluscolor}{HTML}{EEEEEE}
\definecolor{sectioncolor}{HTML}{B2B2B2}
\definecolor{subsectioncolor}{HTML}{B2B2B2}
\definecolor{engcolor}{HTML}{EEEEEE}
\definecolor{exocolor}{HTML}{3B3838}
\definecolor{exosoltitlecolor}{rgb}{0,0,0}
\definecolor{titlecolor}{rgb}{0,0,0}
}




%Exercices résolus dans le cours
%#1 : énoncé
%#2 : solution
\newcounter{exosol}
\newcommand{\exosol}[2]{
\stepcounter{exosol}
\begin{tikzpicture}[node distance=0 cm]
\node[fill=backilluscolor,rounded corners=2pt,anchor=south west] (illus) at (0,-0.02)
{\it \textbf{\textcolor{exosoltitlecolor}{Exercice résolu \arabic{exosol}~:~}}};
\node[fill=backilluscolor,rounded corners=2pt,anchor=north west]at(0,0)
{\parbox{\columnwidth-10pt}{#1\par\medskip{\it \textbf{\textcolor{exosoltitlecolor}{Solution~:~}}}#2 }};
\end{tikzpicture}
\bigskip
}

\newcommand{\suite}[1]{
\begin{tikzpicture}[node distance=0 cm]
\node[fill=backilluscolor,rounded corners=2pt,anchor=north west]at(0,0)
{\parbox{\columnwidth-10pt}{{\it \textbf{\textcolor{exosoltitlecolor}{Suite de la solution~:}}}\par#1}};
\end{tikzpicture}
\bigskip
}



%%%%%%%%%%%%%%%%%%%%%%%%%%%%%%%%%%%%%%%%%%%%%%%%%%%%%%%%%%%%%%%%%%%%%%%%%%%%%%%
%Encadrés pour Propriétés, Théorème, Définitions, exemples, exercices

\usepackage{environ}

%\newcounter{propr}
%\newenvironment{propr}{\refstepcounter{propr}\begin{bclogo}[%noborder=true , 
%couleur = white , arrondi = 0.1 ,logo = \bccrayon , barre = snake , tailleOndu = 1.5 , ombre = true , epOmbre = 0.15 , couleurOmbre = black!30]{Propriété \thepropr}}{\end{bclogo}}

\NewEnviron{Prop}[1][]{
\begin{tikzpicture}[node distance=0 cm]
\node[fill=theocolor,rounded corners=5pt,anchor=south west] (theorem) at (0,0)
{\textcolor{titlecolor}{Propriété~:~#1}};
\node[draw,drop shadow,color=theocolor,very thick,fill=backtheocolor,rounded corners=5pt,anchor=north west] at(0,-0.02)
{\black\parbox{\columnwidth-12pt}{\BODY}};
\end{tikzpicture}
\bigskip
}

%\newcommand{\propr}[2]{
%\begin{tikzpicture}[node distance=0 cm]
%\node[fill=theocolor,rounded corners=5pt,anchor=south west] (theorem) at (0,0)
%{\textcolor{titlecolor}{Propriété~:~#1}};
%\node[draw,drop shadow,color=theocolor,very thick,fill=backtheocolor,rounded corners=5pt,anchor=north west] at(0,-0.02)
%{\black\parbox{\columnwidth-12pt}{#2}};
%\end{tikzpicture}
%\bigskip
%}

%\newenvironment{propr}[1][]{\begin{bclogo}[noborder=true , 
%couleur = white , arrondi = 0.1 , barre = none , ombre = true , epOmbre = 0.15 , couleurOmbre = black!30]{Propriété {\bf \xspace : #1}}}{\end{bclogo}}
%
%
%\newcounter{theo}
%\newenvironment{theo}{\refstepcounter{theo}\begin{bclogo}[%noborder=true , 
%couleur = white , arrondi = 0.1 ,logo = \bcbook , barre = snake , tailleOndu = 1.5 , ombre = true , epOmbre = 0.15 , couleurOmbre = orange]{Théorème \thetheo}}{\end{bclogo}}

%Théorème
\NewEnviron{Thm}[1][]{
\begin{tikzpicture}[node distance=0 cm]
\node[fill=theocolor,rounded corners=5pt,anchor=south west] (theorem) at (0,0)
{\textcolor{titlecolor}{Théorème~:~#1}};
\node[draw,drop shadow,color=theocolor,very thick,fill=backtheocolor,rounded corners=5pt,anchor=north west] at(0,-0.02)
{\black\parbox{\columnwidth-12pt}{\BODY}};
\end{tikzpicture}
\bigskip
}



%\newcounter{defi}
%\newenvironment{defi}{\refstepcounter{defi}\begin{bclogo}[%noborder=true , 
%couleur = white , arrondi = 0.1 ,logo = \bcinfo , barre = snake , tailleOndu = 1.5 , ombre = true , epOmbre = 0.15 , couleurOmbre = blue]{Définition \thedefi}}{\end{bclogo}}
\NewEnviron{Defi}[1][]{
\begin{tikzpicture}[node distance=0 cm]
\node[fill=theocolor,rounded corners=5pt,anchor=south west] (theorem) at (0,0)
{\textcolor{titlecolor}{Définition~:~#1}};
\node[draw,drop shadow,color=deficolor,very thick,fill=backdeficolor,rounded corners=5pt,anchor=north west] at(0,-0.02)
{\black\parbox{\columnwidth-12pt}{\BODY}};
\end{tikzpicture}
\bigskip
}

%\newcounter{methode}
%\newenvironment{methode}{\refstepcounter{methode}\begin{bclogo}[%noborder=true , 
%couleur = white , arrondi = 0.1 ,logo = \bccrayon , barre = snake , tailleOndu = 1.5 , ombre = true , epOmbre = 0.15 , couleurOmbre = red]{Méthode \themethode}}{\end{bclogo}}
\NewEnviron{Methode}[1][]{
\begin{tikzpicture}[node distance=0 cm]
\node[fill=theocolor,rounded corners=5pt,anchor=south west] (theorem) at (0,0)
{\textcolor{titlecolor}{Méthode~:~#1}};
\node[draw,drop shadow,color=methcolor,very thick,fill=backmethcolor,rounded corners=5pt,anchor=north west] at(0,-0.02)
{\black\parbox{\columnwidth-12pt}{\BODY}};
\end{tikzpicture}
\bigskip
}


%chapitres
\makeatletter

\renewcommand{\@makechapterhead}[1]{
\begin{tikzpicture}
\node[fill=theocolor,rectangle,rounded corners=5pt]{%
\begin{minipage}{\linewidth}
\begin{center}
\vspace*{9pt}
\textcolor{titlecolor}{\Large \textsc{\textbf{Chapitre \thechapter \ : \ #1}}}
\vspace*{9pt}
\end{center}
\end{minipage}
};\end{tikzpicture}
}

\makeatother

\NewEnviron{Exemple}[1][]{
\begin{tikzpicture}[node distance=0 cm]
\node[draw,drop shadow,color=methcolor,very thick,fill=backmethcolor,rounded corners=5pt,anchor=north west] at(0,-0.02)
{\black\parbox{\columnwidth-12pt}{\textbf{Exemple~:~#1}\\
\BODY}};
\end{tikzpicture}
\bigskip
}

\newcounter{exemple}\newcommand{\exemple}{\refstepcounter{exemple}\textbf{Exemple \theexemple \ :}\xspace}
\newcounter{exercice}\newcommand{\exercice}{\refstepcounter{exercice}\textbf{\large{Exercice \theexercice \ :}}\xspace}

\newcounter{probleme}\newcommand{\probleme}{\refstepcounter{probleme}\textbf{\large{Problème \theprobleme \ :}}\xspace}
\newcounter{remarque}\newcommand{\remarque}{\refstepcounter{remarque}\textbf{Remarque \theremarque \ :}\xspace}

\NewEnviron{Rmq}{
\textbf{\large{Remarque :}}\par
\BODY
\bigskip
}

\newcounter{rem}\newcommand{\rem}{\refstepcounter{rem}\textbf{R \therem \ :}\xspace}

%Exercices sans numérotation automatique
%\newcommand{\Exercice}[1]{\textbf{\large{Exercice #1 :}\xspace}}
%Exercices du contrôle numérotés
\newcounter{exercice}
\NewEnviron{Exercice}{
\refstepcounter{exercice}\textbf{\large{Exercice \theexercice \ :}}\par
\BODY
\bigskip
}

%Exercices non numérotés
\NewEnviron{Exo}[1][]{
\textbf{\large{Exercice #1 \ :}}\par
\BODY
\bigskip
}


%\Leftrightarrow
\newcommand{\Lr}{\Leftrightarrow}

\pagecolor{white}

%Sommaire dans les chapitres
\usepackage{minitoc}

%%%%%%%%%%%%%%%%%%%%%%%%%%%%%%%%%%%%%%%%%%%%%%%%%%%%%%%%%%%%%%%%%%%%%%%%%%%%%%%

%%%%%%%%%%%%%%%%%%%%%%%%%%%%%%%%%%%%%%%%%%%%%%%%%%%%%%%%%%%%%%%%%%%%%%%%%%%%%%%

%entête classique

\newcommand{\tete}{
	{\noindent
      \textbf{BAT 2}     	            %% Auteur
      \hfill \textbf{Année 2013-2014}   \\  	            	%% Année
	}
	%{\noindent
  %    \url{tremulot.maths@gmail.com}                 	%% Module
  %    \hfill \textbf{\today}   \\  % Numero de la feuille
	%}
	\hrule
	\begin{center} 
	\Large \textsc{\textbf{Chapitre 1 : Nombres complexes}}
	\end{center}
	\hrule
	\noindent 
\medskip}
%%%%%%%%%%%%%%%%%%%%%%%%%%%%%%%%%%%%%%%%%%%%%%%%%%%%%%%%%%%%%%%%%%%%%%%%%%%%%%%

%entête contrôle

\fancypagestyle{garde}{% 
\fancyhead[C]{\small\textbf{Nom,  Prénom : }\makebox[5cm]{\dotfill}\hfill \small \textbf{\seconde 1 - 2013-2014}}
\renewcommand{\headrulewidth}{0cm}}

%Pour utiliser cette présentation, décommenter la ligne
%\thispagestyle{garde}
%située après le \begin{document}


\newcommand{\controle}{
\Ovalbox{
\begin{tabularx}{\linewidth}{|>\centering X|>\centering X}
\hline 
\textbf{\large CONTR\^OLE DE MATH\'EMATIQUES N°1} \tabularnewline 
\textit{Vendredi 04 octobre} \tabularnewline
\hline
\end{tabularx}}

\begin{center}
\begin{minipage}{0.8\linewidth}
\begin{center}
\textit{Les calculatrices ne sont pas autorisées.\\
La qualité de la rédaction, la clarté et la précision des raisonnements seront pris en compte dans l'appréciation de la copie.
Le barème est indicatif.}
\end{center}
\end{minipage}
\end{center}

\noindent
%\vspace{1cm}
}

%%%%%%%%%%%%%%%%%%%%%%%%%%%%%%%%%%%%%%%%%%%%%%%%%%%%%%%%%%%%%%%%%%%%%%%%%%%%%%%

%entête DM

\newcommand{\dm}{
\begin{minipage}{0.9\linewidth}
\begin{center}
\begin{tabularx}{\textwidth}{|c|X|c|}
%{@{\extracolsep{\fill}} |l c r|}
		\hline
		BAT 2 & \centering Devoir maison de mathématiques n° 1 & Année 2013/2014 \\
		\hline
		\multicolumn{3}{|c|}{\`A rendre, sur feuille, au plus tard le 15/05} \\
		\hline
		\multicolumn{3}{|l|}{NOM, Prénom : \dotfill} \\
		\hline
\end{tabularx}
\end{center}
\end{minipage}
\hfill
\begin{minipage}{0.09\linewidth}
\begin{pspicture}(2,2)
	\psbarcode{http://www.tremulotmaths.fr}{eclevel=L}{qrcode}
\end{pspicture}
\end{minipage}

\begin{center}
\textbf{Vous pouvez poser vos questions dans les commentaires de l'\href{http://www.tremulotmaths.fr/spip.php?article308}{article sur le site}.}
\end{center}
\noindent
}
%%%%%%%%%%%%%%%%%%%%%%%%%%%%%%%%%%%%%%%%%%%%%%%%%%%%%%%%%%%%%%%%%%%%%%%%%%%%%%%%%%%%%%
%entête bac blanc
\newcommand{\bacblanc}{
\thispagestyle{empty}
\thiswatermark {
\begin{minipage}{0.95\linewidth}
\vspace{25cm}
\begin{center}
\rotatebox{55}{\scalebox{8}{\color[gray]{0.7}OBLIGATOIRE}}
\end{center}
\end{minipage}
}
Lycée Jean-Pierre Timbaud \hfill Février 2014

.\dotfill
\vspace{2cm}
\begin{center}
\textbf{\LARGE{SESSION 2014}}\\
\vspace{1cm}
\textbf{\LARGE{BACCALAUR\'EAT G\'EN\'ERAL BLANC}}\\
\vspace{1cm}
\textbf{\LARGE{Série SCIENTIFIQUE}}\\
\vspace{1cm}
\textbf{\LARGE{MATH\'EMATIQUES}}

\end{center}

\vspace{1cm}

Durée : 4 heures 
\hfill
Coefficient 7

\vspace{1cm}

\dotfill
\begin{center}
L'usage des calculatrices est autorisé pour cette épreuve.\\
\end{center}
\dotfill

\vspace{1cm}


\begin{center}
\begin{bf}

Le candidat doit traiter les quatre exercices.\\
Le candidat est invité à faire figurer sur la copie toute trace de recherche, même incomplète ou non fructueuse, qu'il aura développée.\\
\vspace{12pt}
Il est rappelé que la qualité de la rédaction, la clarté et la précision des raisonnements entreront pour une part importante dans l'appréciation des copies.\\
\vspace{12pt}
Dès que le sujet vous est remis, assurez-vous qu'il est complet, que toutes les pages sont imprimées.\\
Le sujet comporte 5 pages.\\

\end{bf}
\end{center}
\newpage
}
%%%%%%%%%%%%%%%%%%%%%%%%%%%%%%%%%%%%%%%%%%%%%%%%%%%%%%%%%%%%%%%%%%%%%%%%%%%%%%%

%%%%%%%%%%%%%%%%%%%%%%%
%% DEBUT DU DOCUMENT %%
%%%%%%%%%%%%%%%%%%%%%%%

\begin{document}

\setlength\parindent{0mm}
%\tete 		%entête classique
\controle		%entête contrôle
\thispagestyle{garde} %à décommenter pour avoir l'entête des contrôles
%\dm 		%entête DM
%\enlargethispage{2cm}
%\bacblanc 		%entête bac blanc

\renewcommand \footrulewidth{0.2pt}%
\renewcommand \headrulewidth{0pt}%
\pagestyle{fancy}
\fancyhf{}
\pieddepage{BAT 2 - Nombres complexes}{\thepage / \pageref{LastPage}}{Année 2013-2014}

%\thispagestyle{empty} %annuler le style pour la page en cours, soit la première page, ici


%\title{BACCALAUR\'EAT BLANC}
%\date{}%supprime la date si on n'entre rien entre les accolades

%\maketitle %impose le style plain à la page
%\hrule

%%%%%%%%%%%%%%%%%%%%%%%Page 1%%%%%%%%%%%%%%%%%%%%%%%%%%%%%%%%%%%%%%%%
\begin{spacing}{1.2}

%Modèle propriété :
%\fbox{%
%   \begin{minipage}{\textwidth}
%   \textbf{\ul{Définition :}} \\
%   L'ensemble des matrices d'ordre $(n,p)$ à coefficients réels se note $\mathcal{M}_{n,p}(\R)$.
%   \end{minipage}%
%}


%Autre possibilité avec les packages de la partie ENCADRES
%\begin{bclogo}[couleur = gray!30 , arrondi = 0.1 ,logo = \bclampe , barre = snake , tailleOndu = 1.5]{Théorème}
%Soit $M$ la matrice de transition d'un graphe probabiliste.
%Soit $P_0$ la matrice ligne décrivant l'état initial.
%La matrice ligne $P_n$ décrivant l'état probabiliste à l'étape $n$ vérifie :
%$\bullet P_{n+1}=P_n\times M$
%$\bullet P_n=P_0\times M^n$
%\end{bclogo}
%%%%%%%%%%%%%%%%%%%%%%%%%%%%%%%%%%%%%%%%%%%%%%%%%%%%%%%%%%%%
\renewcommand{\contentsname}{Sommaire}\tableofcontents \newpage
%%%%%%%%%%%%%%%%%%%%%%%%%%%%%%%%%%%%%%%%%%%%%%%%%%%%%%%%%%%%
\thiswatermark {
\begin{minipage}{0.95\linewidth}
\vspace{25cm}
\begin{center}
\rotatebox{55}{\scalebox{8}{\color[gray]{0.7}\LaTeX 2}}
\end{center}
\end{minipage}
}

\section{Introduction}

\subsection{Identification de la situation}

\begin{center}
\psset{xunit=0.5}
\begin{pspicture}(-5,-0.5)(5,0.5)
\psline[linecolor=red,linewidth=3pt]{]-]}(-2.5,0)(3,0)
\psaxes[yAxis=false,arrows=->](0,0)(-5,0)(5,0)
\end{pspicture}
\end{center}


\begin{center}
\begin{pspicture}(2,2)
	\psbarcode{http://www.tremulotmaths.fr}{eclevel=M}{qrcode}
\end{pspicture}

\end{center}
On se situe, dans ce chapitre, dans deux domaines des statistiques, qui sont ceux de \og l'échantillonnage\fg~ et de \og l'estimation\fg~.

Une ligne de pointillés :

\Pointilles

Trois lignes de pointillés :

\Pointilles[3]

\begin{minipage}{0.45\linewidth}

Une ligne de pointillés dans minipage :
\Pointilles

\end{minipage}
\hfill
\begin{minipage}{0.45\linewidth}

Trois lignes de pointillés dans minipage :
\Pointilles[3]

\end{minipage}

Apprenons à reconnaître les différents contextes d'applications.

\medskip


\textbf{\textit{Exemples :}}

On considère deux urnes $U_1$ et $U_2$ contenant chacune un très grand nombre de boules, rouges ou bleues.

\begin{enumerate}
\item 
Dans l'urne $U_1$, \textbf{on connaît la proportion $p$ de boules rouges}.

\begin{minipage}{0.45\linewidth}
On procède à des tirages avec remise de $n$ boules, et on 

observe la fréquence d'apparition d'une boule rouge.

Cette fréquence observée appartient \og en général\fg~ à un \og intervalle de fluctuation \fg~ de centre $p$, dont la longueur diminue lorsque $n$ augmente.

\end{minipage}
\hfill
\begin{minipage}{0.45\linewidth}
\includegraphics[width=5cm]{fig.eps}

\end{minipage}


Cet intervalle est un \og intervalle de fluctuation \fg~.

On est ici dans le domaine de \textbf{l'échantillonnage} et de l'\textbf{intervalle de fluctuation}.
\[\intervalleff{\dfrac{5}{2}}{\sqrt{\dfrac{1}{3}}}\]
\[\intervallefo{\dfrac{5}{2}}{\sqrt{\dfrac{1}{3}}}\]
\[\intervalleof{\dfrac{5}{2}}{\sqrt{\dfrac{1}{3}}}\]
\[\intervalleoo{\dfrac{5}{2}}{\sqrt{\dfrac{1}{3}}}\]

\begin{center}
\definecolor{cqcqcq}{rgb}{0.75,0.75,0.75}
\begin{tikzpicture}[line cap=round,line join=round,>=triangle 45,x=1.0cm,y=1.0cm]
\draw [color=cqcqcq,dash pattern=on 2pt off 2pt, xstep=1.0cm,ystep=1.0cm] (-0.46,-3.28) grid (6.88,3.36);
\draw[->,color=black] (-0.46,0) -- (6.88,0);
\foreach \x in {,1,2,3,4,5,6}
\draw[shift={(\x,0)},color=black] (0pt,2pt) -- (0pt,-2pt) node[below] {\footnotesize $\x$};
\draw[->,color=black] (0,-3.28) -- (0,3.36);
\foreach \y in {-3,-2,-1,1,2,3}
\draw[shift={(0,\y)},color=black] (2pt,0pt) -- (-2pt,0pt) node[left] {\footnotesize $\y$};
\draw[color=black] (0pt,-10pt) node[right] {\footnotesize $0$};
\clip(-0.46,-3.28) rectangle (6.88,3.36);
\draw[smooth,samples=100,domain=0.001:6.88] plot(\x,{ln((\x))});
\end{tikzpicture}
\end{center}

\item Dans l'urne $U_2$, \textbf{on ignore la proportion de boules rouges}.

En procédant à des tirages avec rmise de $n$ boules, on va essayer d'estimer la proportion $p$ de boules rouges dans l'urne.

Cette estimation se fait au moyen d'un \og intervalle de confiance \fg~.

Cet intervalle dépend d'un coefficient, le \og niveau de confiance\fg~, que l'on attribue à l'estimation.

On est ici dans le domaine de l'\textbf{estimation}, et de l'\textbf{intervalle de confiance}.

\begin{center}
\definecolor{cqcqcq}{rgb}{0.75,0.75,0.75}
\begin{tikzpicture}[line cap=round,line join=round,>=triangle 45,x=1.0cm,y=1.0cm]
\draw [color=cqcqcq,dash pattern=on 4pt off 4pt, xstep=2.0cm,ystep=2.0cm] (-1.95,-1.49) grid (8.29,6.99);
\draw[->,color=black] (-1.95,0) -- (8.29,0);
\foreach \x in {,2,4,6,8}
\draw[shift={(\x,0)},color=black] (0pt,2pt) -- (0pt,-2pt) node[below] {\footnotesize $\x$};
\draw[->,color=black] (0,-1.49) -- (0,6.99);
\foreach \y in {,2,4,6}
\draw[shift={(0,\y)},color=black] (2pt,0pt) -- (-2pt,0pt) node[left] {\footnotesize $\y$};
\draw[color=black] (0pt,-10pt) node[right] {\footnotesize $0$};
\clip(-1.95,-1.49) rectangle (8.29,6.99);
\draw[smooth,samples=100,domain=0.01:8.29] plot(\x,{ln((\x))});
\draw[smooth,samples=100,domain=0:1.95] plot(\x,{2.718^(\x)});
\draw[smooth,samples=100,domain=-1.95:0] plot(\x,{2.718^(\x)});
\draw[smooth,samples=100,domain=-1.95:8.29] plot(\x,{(\x)});
\draw (5.07,3.19) node[anchor=north west] {$y=\ln (x)$};
\draw (4.01,6.02) node[anchor=north west] {$y=x$};
\draw (1.14,5.28) node[anchor=north west] {$y=\text{e}^x$};
\end{tikzpicture}

\psset{xunit=2.0cm,yunit=2.0cm}
\begin{pspicture*}(-0.98,-0.74)(4.14,3.49)
\psgrid[subgriddiv=0,gridlabels=0,gridcolor=lightgray](0,0)(-0.98,-0.74)(4.14,3.49)
\psset{xunit=1.0cm,yunit=1.0cm,algebraic=true,dotstyle=o,dotsize=3pt 0,linewidth=0.8pt,arrowsize=3pt 2,arrowinset=0.25}
\psaxes[%labelFontSize=\caligiptstyle,
xAxis=true,yAxis=true,Dx=2,Dy=2,ticksize=-2pt 0,subticks=2]{->}(0,0)(-1.95,-1.49)(8.29,6.99)
\psplot[plotpoints=200]{1.1299160078328312E-8}{8.289582400000004}{ln(x)}
\psplot[plotpoints=200]{-1.953261199999997}{8.289582400000004}{EXP(x)}
\psplot[plotpoints=200]{-1.953261199999997}{8.289582400000004}{x}
\rput[tl](5.07,3.19){$y=\ln (x)$}
\rput[tl](4.01,6.02){$y=x$}
\rput[tl](1.14,5.28){$y=\text{e}^x$}
\end{pspicture*}

\end{center}

\medskip

\begin{minipage}{0.9\linewidth}
\begin{alterqcm}
\AQquestion{Question}{{Proposition 1},{Proposition 2},{Proposition 3}}
\AQquestion{Question}{{Proposition 1},{Proposition 2},{Proposition 3}}
\AQquestion{Question}{{Proposition 1},{Proposition 2},{Proposition 3}}
\AQquestion{Question}{{Proposition 1},{Proposition 2},{Proposition 3}}
\end{alterqcm}
\end{minipage}


\begin{minipage}{0.95\linewidth}
\begin{alterqcm}[VF,lq=125mm,correction,
symb = \dingsquare,
corsymb = \dingchecksquare]
\AQquestion[br=1]{Pour tout $x \in ]-3~;~2],~f'(x) \geqslant 0$.}
\AQquestion[br=2]{La fonction $F$ présente un maximum en $2$}
\AQquestion[br=2]{$\displaystyle\int_{0}^2 f'(x)\:\text{d}x = - 2$}
\end{alterqcm}
\end{minipage}

\end{enumerate}
%%%%%%%%%%%%%%%%%%%%%%%%%%%%%%%%%%%%%%%%%%%%%%%%%%%%%%%%%%%%%%%%%%%%%%%%%
%%%%%%%%%%%%%%%%%%%%%%%%%%%%%%%%%%%%%%%%%%%%%%%%%%%%%%%%%%%%%%%%%%%%%%%%%
\newpage
\subsection{Intervalle de fluctuation ou intervalle de confiance ?}


\begin{center}
\includegraphics[width=0.7\textwidth]{image}
\end{center}

On s'intéresse à une population, dont on étudie un caractère particulier.

\medskip

\begin{bclogo}[couleur = gray!30 , arrondi = 0.1 ,logo = \bclampe , barre = snake , tailleOndu = 1.5]{Deux cas}
\begin{itemize}
\item[$\bullet$] \textbf{\'Echantillonnage}

On utilise un \textbf{intervalle de fluctuation} quand :
	\begin{itemize}
	\item on \textbf{connaît la proportion} $p$ de présence du caractère étudié dans la population,
	\item OU, on \textbf{fait une hypothèse sur la valeur de cette proportion} (on est alors dans le cas de la \og prise de décision\fg~ .
	\end{itemize}

\item[$\bullet$] \textbf{Estimation}

On utilise un \textbf{intervalle de confiance} quand :

on \textbf{ignore la valeur de la proportion} $p$ de présence du caractère dans la population, et on ne formule pas d'hypothèse sur cette valeur.
\end{itemize}
\end{bclogo}

\bigskip
En appliquant l'agorithme de Dijsktra, on obtient :

\begin{center}
\begin{tabular}{|*{8}{>{\centering}m{1.5cm}|}c|}
 \hline
 A&B&C&D&E&F&G&H&Sommet\\
 \hline
 0&$\infty$&$\infty$&$\infty$&$\infty$&$\infty$&$\infty$&$\infty$&A\\
 \hline
  &0+2\par 2(A)&0+1\par 1(A)&0+4\par 4(A)&$\infty$&$\infty$&$\infty$&$\infty$&C\\
 \hline
  &2(A) & &4(A)&1+5\par 6(C)&$\infty$&$\infty$&$\infty$&B\\
 \hline
  & & &4(A)&6(C)&$\infty$&$\infty$&$\infty$&D\\
 \hline
  & & & &6(C)&4+1\par 5(D) &$\infty$&$\infty$&F\\
 \hline
  & & & &6(C)& &5+4\par 9(F)&5+1\par 6(F)&E\\
 \hline
  & & & & & &6+2\par 8(E)&6(F)&H\\
 \hline
  & & & & & &6+4\par 8(E)& &G\\
 \hline
 \end{tabular}
\end{center}



\[\intervalleff{\frac{5}{2}}{-2}\]


\textbf{\textit{Exemples :}}
\begin{enumerate}
\item On dispose d'une pièce de monnaie.

Comment décider qu'elle est \og équilibrée\fg~ ou pas ?

On va ici faire l'hypothèse que la fréquence d'apparition de \og Pile\fg~, par exemple, est égale à 0,5, et on va tester cette hypothèse.

On est dans une \textbf{situation d'échantillonnage}.

\item Une usine fabrique des fusées de feux d'artifice.

Sur 100 fusées choisies au hasard à l'issue du processus de fabrication et mises à feu, on trouve 12 fusées qui ne fonctionnent pas.

Comment se faire une idée de la proportion de fusées défectueuses dans la production ?

On est dans une \textbf{situation d'estimation} : on n'a, au départ, aucune idée de la valeur de la proportion étudiée dans la production.
\end{enumerate}

\medskip


\begin{bclogo}[couleur = blue!30 , arrondi = 0.1 ,logo = \bclampe , barre = snake , tailleOndu = 1.5]{Exercice :}


\textbf{Texte à trous :}

Un enfant a réalisé $\np{4040}$ \trou{lancers} d'une pièce de monnaie, et il a obtenu $\np{2048}$ fois le résultat \og Pile\fg~.

Nous allons nous intéresser au problème suivant : \og Peut-on considérer que la pièce utilisée est équilibrée ?\fg~

\begin{enumerate}
\item Prendre une décision en utilisant l'intervalle de fluctuation au seuil de 95\% sous la forme étudiée en classe de 2\up{nde}.
\item Reprendre le même travail, en utilisant un intervalle de fluctuation au seuil 95\% obtenu à partir de la loi binomiale vue en classe de 1\up{ère}.
\end{enumerate}
\end{bclogo}


\newpage

\textbf{\textit{Solution :}}

\begin{enumerate}
\item En classe de 2\up{nde}, on a vu que pour $n>25$ et $0,2<p<0,8$, l'intervalle $\left[p-\frac{1}{\sqrt{n}};p+\frac{1}{\sqrt{n}}\right]$ constitue un intervalle de fluctuation au seuil 95\%.

Ici, si on appelle $p$ la \og proportion théorique\fg~ de résultats \og Pile\fg~ obtenus sur un très grand nombre de lancers, on veut tester l'hypothèse selon laquelle la pièce est équilibrée, c'est-à-dire l'hypothèse $p=0,5$.

Pour l'échantillon considéré, on a $n=\np{4040}$ ; les conditions d'application sont donc réalisées, et on obtient l'intervalle de fluctuation 
\[\left[0,5-\frac{1}{\sqrt{\np{4040}}};0,5+\frac{1}{\sqrt{\np{4040}}}\right]\]
En arrondissant à $10^{-4}$ près, on obtient l'intervalle de fluctuation $I_1=[\np{0,4843};\np{0,5157}]$.

La fréquence observée dans l'échantillon est donnée par :
\[f=\frac{\np{2048}}{\np{4040}}\approx \np{0,5069}~\text{à}~10^{-4}~\text{près}\]

Ainsi, on a $f\in I_1$, donc on accepte l'hypothèse de pièce équilibrée au seuil de confiance de 95\%.

\item En supposant que la pièce est équilibrée, la variable aléatoire $X$ qui dénombre les résultat \og Pile\fg~ obtenus parmi les $\np{4040}$ lancers réalisés suit une loi binomiale de paramètres $n=\np{4040}$ et $p=0,5$.

On a vu en classe de 1\up{ère} que l'ntervalle de fluctuation à 95\% associé à la variable aléatoire $X$ est l'intervalle $I_2=\left[\frac{k_1}{n};\frac{k_2}{n}\right]$, où $k_1$ est le plus petit entier $k$ vérifiant $P(X\leq k)> 0,025$ et $k_2$ le plus petit entier $k$ vérifiant $P(X\leq k)\geq 0,975$.

On détermine $k_1$ et $k_2$ à l'aide d'un tableur (ou d'une calculatrice...mais le tableau est très long, ici).

On obtient $k_1=\np{1958}$ et $k_2=\np{2082}$, qui donnent comme intervalle de fluctuation :
\[I_2=\left[\frac{\np{1958}}{\np{4040}};\frac{\np{2082}}{\np{4040}}\right]\]
soit en arrondissant à $10^{-4}$ près $I_2=[\np{0,4847};\np{0,5153}]$.

La fréquence observée $f\approx \np{0,5069}$ vérifie $f\in I_2$.

On est ici aussi conduit à accepter l'hypothèse de pièce équilibrée au seuil de confiance de 95\%.
\end{enumerate}
%%%%%%%%%%%%%%%%%%%%%%%%%%%%%%%%%%%%%%%%%%%%%%%%%%%%%%%%%%%%%%%%%%%%%%%%%

\newpage
%%%%%%%%%%%%%%%%%%%%%%%%%%%%%%%%%%%%%%%%%%%%%%%%%%%%%%%%%%%%
\textbf{\textsc{Exercice 2 :}\hfill 8 points}% ex 109 p 29 indice
%%%%%%%%%%%%%%%%%%%%%%%%%%%%%%%%%%%%%%%%%%%%%%%%%%%%%%%%%%%%
\medskip

Une agence de voyages propose un circuit touristique pour visiter trois villes A, B et C.

Le client peut choisir la durée de son séjour dans chacune des villes.

L'agence propose des tarifs qui différent selon la période. Il existe trois périodes touristiques : une période haute, une période moyenne et une période basse.

Les prix journaliers, en \textbf{centaines d'euros} par personne, dans les différents lieux sont donnés dans le tableau suivant :

\begin{center}
\newcolumntype{K}[1]{>{\columncolor{#1}}l}
\begin{tabular}{|K{gray}|*{3}{c|}}
\cline{2-4}
\multicolumn{1}{c|}{} & \cellcolor{yellow} \bf Ville A & \cellcolor{yellow}\bf Ville B & \cellcolor{yellow}\bf Ville C\\
\hline
\bf Période haute & 2,5 & 3,5 & 1,5\\
\hline
\bf Période moyenne & 2 & 2,2 & 1,5\\
\hline
\bf Période basse & 1,5 & 1,2 & 1\\
\hline
\end{tabular}
\end{center}

On considère la matrice $M=\begin{pmatrix}
2,5&3,5&1,5\\2&2,2&1,5\\1,5&1,2&1
\end{pmatrix}$

\begin{enumerate}
\item M. Gibbs choisit un circuit de trois jours : deux jours dans la ville A et un jour dans la ville C.

Calculer, à l'aide d'un \textbf{produit de matrices}, le prix de son séjour pour chaque période.

\item La matrice ci-dessous donne toutes les combinaisons d'un circuit de trois jours :

\[D=\begin{pmatrix}
3&2&2&1&1&1&0&0&0&0\\
0&1&0&2&0&...&...&...&...&...\\
0&0&1&0&2&...&...&...&...&...
\end{pmatrix}\]

Ainsi, la troisième colonne indique que M. Gibbs choisit un circuit de deux jours dans la ville A et un jour dans la ville C.

	\begin{enumerate}
	\item Compléter la matrice $D$.
	
	\item Calculer la matrice $M\times D$.
	
	\end{enumerate}
\item En déduire quels sont les circuits de trois jours possibles pour que le séjour coûte moins de :
	\begin{enumerate}
	\item 700 \euro~ en période haute ;
	
	\item 560 \euro~ en période moyenne ;
	
	\item 330 \euro~ en période basse.
	
	\end{enumerate}

\item Dans une publicité, l'agence de voyages affirme qu'un circuit complet de 10 jours est possible au prix de \np{2800} \euro~ en période haute, \np{2000} \euro~ en période moyenne et \EUR{\np{1250}} en période basse.

Comment se compose ce voyage ?

\end{enumerate}



%%%%%%%%%%%%%%%%%%%%%%%%%%%%%%%%%%%%%%%%%%%%%%%%%%%%%%%%%%%%%%%%%%%%%%%%%
\newpage
\section{Intervalles de fluctuation}

\subsection{Intervalle de fluctuation asymptotique}

On dispose d'un urne contenant un très grand nombre de boules rouges et bleues.

On sait que la proportion de boules rouges dans l'urne est égale à $p=0,4$.

Si on tire, successivement avec remise, $n$ boules dans l'urne ($n$ entier naturel non nul), et si on appelle $X_n$ la variable aléatoire dénombrant les boules rouges tirées, alors $X_n$ suit une loi binomiale $\calig{B}(n;p)$.

\medskip

\begin{theo}
Soit $X_n$ une variable aléatoire suivant une loi binomiale $\calig{B}(n;p)$, et $\alpha$ un réel tel que $0<\alpha<1$.

Si $X$ est une variable aléatoire suivant la loi normale centrée réduite $\calig{N}(0;1)$, on appelle $u_{\alpha}$ l'unique réel tel que :
\[P(-u_{\alpha}\leq X\leq u_{\alpha})=1-\alpha\]
\begin{itemize}
\item[$\bullet$] On appelle $I_n$ l'intervalle :
\[I_n=\left[p-u_{\alpha}\frac{\sqrt{p(1-p)}}{\sqrt{n}};p+u_{\alpha}\frac{\sqrt{p(1-p}}{\sqrt{n}}\right]\]
Alors,
\[\limn P\left(\frac{X_n}{n}\in I_n=1-\alpha\right)\]

\item[$\bullet$] L'intervalle $I_n$ contient la fréquence $F_n=\frac{X_n}{n}$ avec une probabilité qui se rapproche de $1-\alpha$ lorsque $n$ augmente : 

on dit que c'est l'\textbf{intervalle de fluctuation asymptotique de $F_n$ au seuil $1-\alpha$}.
\end{itemize}
\end{theo}

\bigskip

\begin{bclogo}[couleur = yellow!30 , arrondi = 0.1 ,logo = \bclampe , barre = snake , tailleOndu = 1.5]{Démonstration (exigible) :}

On pose $Z_n=\frac{X_n-np}{\sqrt{np(1-p)}}$ et on applique le théorème de Moivre-Laplace.

Si $X$ suit la loi normale $\calig{N}(0;1)$ : $\limn P\left(Z_n\in [-u_{\alpha};u_{\alpha}]\right)=P(X\in [-u_{\alpha};u_{\alpha}]=1-\alpha$.

Or, $Z_n\in [-u_{\alpha};u_{\alpha}] \qLRq -u_{\alpha}\leq \frac{X_n-np}{\sqrt{np(1-p)}}\leq u_{\alpha} \qLRq np-u_{\alpha}\sqrt{np(1-p}\leq X_n \leq np+u_{\alpha}\sqrt{np(1-p)}$.

$Z_n\in [-u_{\alpha};u_{\alpha}] \qLRq p-u_{\alpha}\frac{\sqrt{p(1-p}}{\sqrt{n}}\leq \frac{X_n}{n}\leq p+u_{\alpha}\frac{\sqrt{p(1-p)}}{\sqrt{n}}$.

Donc, $Z_n\in [-u_{\alpha};u_{\alpha}] \qLRq \frac{X_n}{n}\in I_n$.
\end{bclogo}

\bigskip

\textbf{\textit{Exemple :}}

On tire 50 boules de l'urne décrite ci-dessus, et on souhaite déterminer un intervalle de fluctuation au seuil 0,9 (c'est-à-dire $\alpha=0,1$).

À l'aide de la calculatrice, on trouve 1,645 comme valeur approchée à $10^{-3}$ près de $u_{0,1}$.

On obtient donc comme intervalle de fluctuation :
\[I_{50}=\left[0,4-u_{0,1}\frac{\sqrt{0,4\times 0,6}}{\sqrt{50}};0,4+u_{0,1}\frac{\sqrt{0,4\times 0,6}}{\sqrt{50}}\right]\]
soit $I_{50}=[0,286;0,514]$.

Ainsi, en effectuant 50 tirages dans cette urne, la fréquence d'apparition d'une boule rouge est comprise entre 0,28 et 0,514 avec une probabilité d'environ 0,9.

Pour 500 tirages, au même seuil de 0,9, on obtient $I_{500}=[0,364;0,436]$.

\medskip

\textbf{\textit{Remarque :}}

La longueur de l'intervalle a, pour un même seuil, été divisée par plus de 3 en passant de 50 à 500 tirages \quad !!

\medskip

\begin{propr}

L'intervalle de fluctuation asymptotique au seuil de 95\% pour une variable aléatoire $X_n$ suivant une loi binomiale $\calig{B}(n;p)$ est l'intervalle :
\[I_n=\left[p-1,96\frac{\sqrt{p(1-p)}}{\sqrt{n}};p+1,96\frac{\sqrt{p(1-p)}}{\sqrt{n}}\right]\]

\end{propr}

\bigskip

\textbf{\textit{Démonstration :}}

Cette propriété découle de celle vue au chapitre 11 concernant la loi normale.

On a vu en effet que $u_{0,05}\approx 1,96$.

\medskip

\textbf{\textit{Remarque :}}

En majorant $1,96\sqrt{p(1-p)}$ par 1, on retrouve l'intervalle de fluctuation présenté en classe de 2\up{nde}.

La majoration précédente résulta de ce qui suit :

$p(1-p)=-p^2+p$ est un trinôme du second degré.

Il admet un maximum en $p=0,5$, qui est égal à 0,25.

Ainsi, pour tout $p\in [0;1], 0\leq p(1-p)\leq 0,25$.

La fonction racine carrée étant croissante sur $\R_+$, on en déduit : $0\leq \sqrt{p(1-p)}\leq 0,5$, 

puis $0\leq 1,96\sqrt{p(1-p)}\leq 1,95\times 0,5\leq 2\times 0,5=1$.

%%%%%%%%%%%%%%%%%%%%%%%%%%%%%%%%%%%%%%%%%%%%%%%%%%%%%%%%%%%%%%%%%%%%%%%%%
%%%%%%%%%%%%%%%%%%%%%%%%%%%%%%%%%%%%%%%%%%%%%%%%%%%%%%%%%%%%%%%%%%%%%%%%%
\subsection{Exercice}

\begin{defi}

Dans un casino, il a été décidé que les \og machines à sous\fg~ doivent être réglées sur une fréquence de gain du joueur de $g=0,06$.

Une fréquence inférieure est supposée faire \og fuir \fg~ le client, et une fréquence supérieure est susceptible de ruiner le casino.

Trois contrôleurs différents vérifient une même machine.

Le premier a joué 50 fois et gagné 2 fois,

le second a joué 120 fois et gagné 14 fois,

le troisième a joué 400 fois et gagné 30 fois.

En utilisant des intervalles de fluctuation asymptotiques au seuil 95\%, examiner dans chaque cas la décision à prendre par le contrôleur, à savoir accepter ou rejeter l'hypothèse $g=0,06$.
\end{defi}

\bigskip

\textbf{\textit{Solution :}}

\begin{itemize}
\item[$\bullet$] \textbf{Premier contrôleur}
Le contrôle a porté sur 50 parties, donc $n=50$.

On a bien $n\geq 30$, mais $np=50\times 0,06=3$ et donc $np<5$.

On n'est pas dans les conditions d'application d'un intervalle de fluctuation asymptotique.

Ce contrôle ne peut rien donner de probant en termes de prise de décision.

\item[$\bullet$] \textbf{Deuxième contrôleur}

Le contrôle a porté sur 120 parties, donc $n=120, np=7,2$ et $n(1-p)=112,8$.

Les conditions d'utilisation d'un intervalle de fluctuation asymptotique sont donc réunies.

L'intervalle de fluctuation asymptotique au seuil 95\% pour un échantillon de taille $n=120$ est égal à :
\[I_{120}=\left[0,06-1,96\frac{\sqrt{0,06\times 0,94}}{\sqrt{120}};0,06+1,96\frac{\sqrt{0,06\times 0,94}}{\sqrt{120}}\right]\]
soit en arrondissant les bornes à $10^{-4}$ près : $I_{120}=[\np{0,0175};\np{0,1025}]$.

La fréquence observée par le deuxième contrôleur est $f_2=\frac{14}{120}\approx \np{0,1167}$.

Cette fréquence est à l'extérieur de l'intervalle de fluctuation $I_{120}$, et le second contrôleur est conduit à rejeter l'hypothèse $g=0,06$.

\item[$\bullet$] \textbf{Troisième contrôleur}

La contrôle a porté sur 400 parties, donc $n=400, np=24$ et $n(1-p)=376$.

Les conditions d'utilisation d'un intervalle de fluctuation asymptotique sont donc réunies.

L'intervalle de fluctuation asymptotique au seuil 95\% pour un échantillon de taille $n=400$ est égal à :

\[I_{400}=\left[0,06-1,96\frac{\sqrt{0,06\times 0,94}}{\sqrt{400}};0,06+1,96\frac{\sqrt{0,06\times 0,94}}{\sqrt{400}}\right]\]
soit en arrondissant les bornes à $10^{-4}$ près : $I_{400}=[\np{0,0367};\np{0,0833}]$.

La fréquence observée par le deuxième contrôleur est $f_3=\frac{30}{400}= \np{0,075}\in I_{400}$.

Ce dernier est donc amené à accepter l'hypothèse $g=0,06$.

\end{itemize}

\medskip

\begin{bclogo}[couleur = gray!30 , arrondi = 0.1 ,logo = \bclampe , barre = snake , tailleOndu = 1.5]{Remarques}

\begin{enumerate}
\item Pour tester une hypothèse à l'aide d'un intervalle de fluctuation asymptotique, on doit vérifier ses conditions d'utilisation :
\[n\geq 30, \quad np\geq 5, \quad n(1-p)\geq 5\]

Ces conditions vérifiées, on détermine l'intervalle de fluctuation $I_n$ au seuil demandé (souvent 95\%).

Si la fréquence observée $f$ n'appartient pas à $I_n$, on rejette l'hypothèse.

\item L'intervalle de fluctuation étudié est \og bilatéral\fg~ : on est conduit à rejeter l'hypothèse dans le cas où la fréquence observée est \og trop grande\fg~, ou \og trop petite \fg~.
\end{enumerate}
\end{bclogo}

%%%%%%%%%%%%%%%%%%%%%%%%%%%%%%%%%%%%%%%%%%%%%%%%%%%%%%%%%%%%%%%%%%%%%%%%%


%%%%%%%%%%%%%%%%%%%%%%%%%%%%%%%%%%%%%%%%%%%%%%%%%%%%%%%%%%%%

\end{spacing}

%%%%%%%%%%%%%%%%%%%%%%%%%%%%%%%%%%%%%%%%%%%%%%%%%%%%%%%%%%%%
%%%%%%%%%%%%%%%%%%%%%
%% FIN DU DOCUMENT %%
%%%%%%%%%%%%%%%%%%%%%
\end{document}